\chapter{Interfaces}

CL provides many mechanisms for interfacing with the ``outside world.'' This chapter will
introduce some of the more important and practical ones available as part of Allegro CL\index{Allegro CL},
or as extensions to Allegro CL through third parties. These techniques include:

\begin{itemize}
\item Direct operating system calls
\item Foreign Function Interface (FFI) for calling functions written in other languages
\item CORBA\index{CORBA}
\item TCP/IP Sockets
\item Database interfaces
\item Web Serving
\end{itemize}

\section{The Operating System}

\subsection{System Calls}

CL provides the function \texttt{excl:run-shell-command} for making direct calls to the operating
system. Note that the \texttt{excl:} package qualifier is used, meaning that this call is an extension
to Common Lisp\index{Common Lisp} and is specific to Allegro CL. Other CL platforms provide similar functionality. Calling OS commands directly is generally a convenient way of interfacing with the
Unix (Linux) operating system or with legacy applications. It is also a quick way to accomplish tasks such as file copying, moving, and deletion without using CL equivalents such
as \texttt{rename-file} and \texttt{delete-file}.

\texttt{Run-shell-command} takes a string representing the desired command, and any optional
keyword arguments as specified in the Allegro CL documentation.

\begin{verbatim}
CL-USER(2): (excl:run-shell-command "ls -l")
total 24
-rwxr-xr-x  1 dcooper users 2863 May  9 14:06 compile-system.cl
-rwxr-xr-x  1 dcooper users 1856 May  8 18:13 example.cl
-rw-r--r--  1 dcooper users  269 May  8 17:48 load-system.cl
-rw-r--r--  1 dcooper users 2152 May  8 18:10 readme.txt
-rwxr-xr-x  1 dcooper users  356 May  8 18:10 start.sh
0
CL-USER(3): (excl:run-shell-command "cp readme.txt readme.bak")
0
\end{verbatim}

\subsection{Environment Variables}

To access values from Unix environment variables, you can use the function \texttt{sys:getenv},
which takes a string representing the name of the environment variable, and returns a string
representing the current value (or NIL if the variable does not exist):

\begin{verbatim}
CL-USER(4): (sys:getenv "SHELL")
"/bin/bash"
\end{verbatim}

You can also set the value of environment variables using \texttt{setf}\index{setf@\texttt{setf}}:

\begin{verbatim}
CL-USER(5): (setf (sys:getenv "MYVAR") "yes")
"yes"
CL-USER(6): (sys:getenv "MYVAR")
"yes"
\end{verbatim}

\section{Foreign Function Interface}

CL provides the capability to call functions written in other languages, most notably the C
language. Each CL platform provides its own Foreign Function Interface (``FFI''), and the
details of this are unfortunately not standardized between platforms. To learn the details
of FFI for Allegro CL, you should consult Allegro CL's documentation.

In general, the steps involved in calling foreign functions from CL are:

\begin{enumerate}
\item Define the Lisp names and types for the C functions and any supporting data structures

\item Load the shared library (\texttt{.so} file on Unix/Linux, or \texttt{.dll} on Windows) which contains
the compiled object code for the functions

\item Call the functions
\end{enumerate}

The FFI capability is extremely useful for ``wrapping'' existing libraries of C code, and/or
for giving CL applications convenient access to low-level system calls. The Franz Inc. website
contains documentation for a large library of Unix system calls which are pre-wrapped and
available from CL.

\section{CORBA}

CORBA\footnote{CORBA stands for Common Object Request Broker Architecture} is a standard for writing distributed object-oriented applications which
can communicate over a network. A CORBA-wrapped application can be written in any
language which supports CORBA. These applications can then communicate with each
other, regardless of the language in which they are written.

CorbaLink is a CORBA implementation for Allegro CL, available through Franz Inc. or
through third parties. Information on CORBA in general can be found at the Object Management Group's website, http://www.omg.org. By using CorbaLink, CL applications can
communicate in a seamless manner with applications written in other CORBA-supporting
languages such as C++, Java, Smalltalk, and others.

\section{TCP/IP Sockets}

TCP/IP provides a standard for communication over a network. Most Internet communication is done via TCP/IP. Both servers and clients can be written which use TCP/IP as the
underlying protocol.

In Allegro CL, the package \texttt{socket} provides CL-level functions which enable the use of
TCP/IP. As a simple example, we will show how to write both a client and a server which
can exchange data using TCP/IP sockets. These two programs can communicate when both
running on the same machine, or they can run on separate machines anywhere on a TCP/IP
network.

First we write the server\footnote{For space considerations we are not handling errors or trying to support multiple connections simultaneously.
See the Allegro CL documentation for socket examples which do provide full error-handling and simultaneous
connection support.}:

\begin{verbatim}
(in-package :user)

(defparameter *socket-server-port* 4444)

(defun open-server ()
  (socket:make-socket :connect :passive
                      :local-port *socket-server-port*))

(defun test-socket-server ()
  (let ((socket (open-server)))
    (unwind-protect
        (let ((stream (socket:accept-connection socket)))
          (let ((command (read stream)))
            (cond ((eql command :hello)
                   (format stream "Hello to you!~%")
                   (finish-output stream))
                  ((eql command :add-1-2)
                   (format stream "~a~%" (+ 1 2))
                   (finish-output stream))
                  (t (format stream "Sorry, I don't understand~%")
                     (finish-output stream))))
          (close stream))
      (close socket))))
\end{verbatim}

Now we write a simple client, which will contact the server, send a command, and print
out the response:

\begin{verbatim}
(in-package :user)

(defparameter *socket-server-port* 4444)

(defun test-socket-client (command)
  (let ((socket (socket:make-socket :remote-host "localhost"
                                    :remote-port *socket-server-port*)))
    (unwind-protect
        (progn
          (format socket "~s~%" command)
          (finish-output socket)
          (format t "~a~%" (read-line socket)))
      (close socket))))
\end{verbatim}

To test the server and client, start two CL processes. Load the server code into one process,
and call

\begin{verbatim}
(test-socket-server)
\end{verbatim}

in that process. Load the client code into the other process, and call

\begin{verbatim}
(test-socket-client :hello)
\end{verbatim}

in that process. In the client, you should see the response from the server printed out. In
the server, you should see the server return from the function call (since the server, as it is
written, only serves one request, then exits). Note that you could change ``localhost'' in the
definition of \texttt{test-socket-client} to any machine name on your network, and the client
and server could communicate with each other between different machines.

Also note that both the \texttt{socket:make-socket} and \texttt{socket:accept-connection} functions accept several optional keyword parameters, which are documented in Allegro CL's
online and hardcopy manuals.

\section{Databases}

Allegro CL provides connectivity to many SQL-based database systems, such as Oracle,
MySQL, PostgreSQL, Sybase and ODBC-compliant databases. Information on these packages, as well as more general information on accessing relational databases from CL, can be
found at

\begin{center}
\texttt{http://opensource.franz.com}
\end{center}

As a short example of database access from CL, we will present code for creating a table,
adding some data to it, and querying the data with MySQL (which is a free open-source
database server available on several platforms including Unix, Linux, and Windows).

\begin{verbatim}
(in-package :user)

(require :mysql)

(use-package :dbi.mysql)

(defun test-mysql ()
  (let ((db (connect "localhost" "testdb" "testuser" "testpasswd")))
    (sql db "create table captable (id int, name char(30))")
    (sql db "insert into captable values (~d, '~a')" 1 "Chicago")
    (sql db "insert into captable values (~d, '~a')" 2 "Austin")
    (sql db "insert into captable values (~d, '~a')" 3 "Philadelphia")
    (dolist (row (sql db "select * from captable where id > 1" :types :auto))
      (format t "~a: ~a~%" (first row) (second row)))
    (sql db "drop table captable")
    (disconnect db)))
\end{verbatim}

In the above example, ``testdb'' is the name of a database which must exist on the MySQL
server running on ``localhost,'' and ``testuser'' is a MySQL username which must have the
proper access privileges to the database (``testpasswd'' is the password).

Upon calling \texttt{(test-mysql)}, we expect to see the following output:

\begin{verbatim}
2: Austin
3: Philadelphia
\end{verbatim}

\section{Webserving}

Most CL vendors either provide direct support or provide third-party support for serving web
content through their CL environments. For Allegro CL, Franz Inc. provides AllegroServe
(``aserve''), a full-featured web application server. Applications built with aserve can provide
all the benefits of traditional web applications, with the additional advantage that all state
can be maintained in the running CL session. For traditional web applications, a new process
must typically be started for each connection; this leads to an inherent inefficiency since all
the application state must be fetched from a database or from persistent files every time.

AllegroServe can be obtained from

\begin{center}
\texttt{http://opensource.franz.com}
\end{center}

and is distributed under an open-source (LLGPL\footnote{LLGPL is the Lisp Lesser General Public License}) license.

\subsection{Publishing Static Pages}

To start the webserver, you use the function \texttt{net.aserve:start}. It takes several keyword
arguments. The only argument required is \texttt{:port}, which specifies the port number that
the server should listen on. The default HTTP port is 80. If you are running as a normal
(non-root) user on Unix/Linux, you will probably not have permission to run on port 80, so
you should choose a number above 1024 and note it for use in forming your URLs. ``8000''
is a traditional ``high port'' default for running webservers. The URL for your webserver
will then appear as \texttt{http://localhost:8000/} (of course replacing 8000 with whatever
port number you actually use).

In order to serve a simple plain HTML (i.e. static) page from a file, the function is
\texttt{publish-file}:

\begin{verbatim}
(net.aserve:start :port 8000)

(publish-file :path "/welcome.html" :file "/tmp/welcome.html")
\end{verbatim}

The \texttt{:path} argument gives the URL path, while the \texttt{:file} argument gives the name in
the Unix (or Windows) filesystem where the file actually resides.

To test the above example, create a simple HTML file and place it in \texttt{/tmp/welcome.html}
(or use the \texttt{:file} argument with a different location as appropriate). A simple HTML file
might look like this:

\begin{verbatim}
<html>
  <head>
    <title>Welcome</title>
  </head>
  <body>
    <h2>Welcome to my Web Page!</h2>
  </body>
</html>
\end{verbatim}

Now, start your web browser and visit the page \texttt{http://localhost:8000/welcome.html}.
You should see the web page come up in the browser.

\subsection{Publishing Dynamic Pages}

Dynamic web page serving is more interesting, since a function is called to generate the
HTML content on-the-fly. This function can tap into the entire running CL environment,
to query data structures, call external databases, etc. Here is an example:

\begin{verbatim}
(publish :path "/hello"
         :content-type "text/html"
         :function
         #'(lambda(req ent)
             (with-http-response (req ent)
               (with-http-body (req ent)
                 (html (:html
                        (:head (:title "Hello Page"))
                        (:body
                         ((:font :color "red")
                          (:h2 "Hello There!")))))))))
\end{verbatim}

To test the above example, visit the URL \texttt{http://localhost:8000/hello} in your browser.
You should see the heading ``Hello There!'' in red in your browser window.

You should also see in your CL console that some information is being logged on each
visit to the page.

Note the use of the \texttt{html} macro, which works somewhat like the \texttt{format}\index{format@\texttt{format}} function. It
takes a series of HTML ``tags'' represented as lists --- the first element of each list is a
keyword symbol representing the ``name'' of the tag, and the rest of the list are either
strings representing actual content, or more nested lists, or both.

\subsection{Query Parameters}

Query parameters are named values which can be spliced into a URL to pass data into
the page. Within AllegroServe they can be accessed with the \texttt{request-query-value}
function. As an example, consider the following \texttt{publish} form:

\begin{verbatim}
(publish :path "/welcome"
         :content-type "text/html"
         :function
         #'(lambda(req ent)
             (with-http-response (req ent)
               (with-http-body (req ent)
                 (let ((name (or (request-query-value "name" req)
                                 "Pal")))
                   (html (:html
                          (:head (:title "Welcome Page"))
                          (:body
                           ((:font :color "red")
                            (:h2 "Welcome, " (:princ name) "!"))))))))))
\end{verbatim}

To test this, visit \texttt{http://localhost:8000/welcome?name=Joe} in your browser. You should
see the heading ``Welcome, Joe!'' in the browser window.

\subsection{Form Parameters}

Query parameters are typically used for small amounts of data. For larger amounts of data,
you will probably want to use HTML form parameters. These are generally filled in by a
user through an interactive HTML form. To demonstrate the server side of such interaction,
we will use the \texttt{curl} Unix command to simulate the posting of form data to a published
aserve page.

Here is a \texttt{publish} form which can accept two form parameters named ``color'' and
``shape'':

\begin{verbatim}
(publish :path "/form-test"
         :content-type "text/html"
         :function
         #'(lambda(req ent)
             (with-http-response (req ent)
               (with-http-body (req ent)
                 (let ((color (request-query-value "color" req))
                       (shape (request-query-value "shape" req)))
                   (html (:html
                          (:head (:title "Form Test"))
                          (:body
                           (:h2 "You posted:")
                           (:princ (format nil "Color: ~a~%" color))
                           (:br)
                           (:princ (format nil "Shape: ~a" shape))))))))))
\end{verbatim}

Now, from a Unix terminal window, type the following command:

\begin{verbatim}
curl --data "color=red&shape=circle" http://localhost:8000/form-test
\end{verbatim}

You should see the following HTML printed on your console:

\begin{verbatim}
<html><head><title>Form Test</title></head>
<body><h2>You posted:</h2>Color: red<br />Shape: circle</body></html>
\end{verbatim}

Note that \texttt{request-query-value} can be used to retrieve both URL-encoded query parameters and form parameters --- aserve makes these available through the same interface.

\subsection{Multipart Form Data and File Uploads}

In order to handle binary files which are uploaded to the server, aserve provides the function
\texttt{get-multipart-content}. Combined with the AllegroServe documentation and/or the
aserve source code, this should be enough to get you started with handling uploaded files.

\subsection{Development and Production Modes}

AllegroServe has a concept of ``development'' and ``production'' mode which can be switched
with the \texttt{:debug} keyword argument. In debug mode, when an error is detected the handler
will stop with a break (see Section 2.8.3). This is usually convenient for development. In
production mode, on the other hand, aserve will simply output a standard HTTP error
response back to the browser and will not stop the server process.

The default debug setting for aserve is production mode. For development, you will
typically want to start aserve with:

\begin{verbatim}
(net.aserve:start :port 8000 :debug t)
\end{verbatim}

\subsection{Other Features}

This chapter has barely scratched the surface of what is available in AllegroServe. For a
complete treatment, see the documentation available at the AllegroServe website listed above.
Additional capabilities of note include:

\begin{itemize}
\item Authorization and security
\item Support for SSL (Secure Sockets Layer)
\item HTTP proxying
\item Virtual hosting (multiple hostnames and DNS names being served by the same physical
server)
\item Support for client-side HTTP from CL
\item Support for internationalization (international character sets) in URLs
\item Support for serving WebDAV (Web-based Distributed Authoring and Versioning)
\end{itemize}

\section{Embedded Languages}

This section applies to layered application development. In particular, we cover the concept
of an Embedded Language: a collection of macros and functions which add primitives and
other key operators to CL, effectively creating a ``language on top of CL.''

An example of this is GDL (Generative Design Language), which is the open-source
core of the Genworks\textregistered{} GDL commercial product. GDL is a general-purpose language
for Knowledge-Based Engineering, and is a significant superset of ANSI CL. Objects in
GDL are similar to CLOS\index{CLOS} objects, but add features such as cached demand-driven behavior,
automatic dependency tracking, and optional default values for unspecified data. GDL adds
new primitives to CL, such as the \texttt{define-object} operator, which is analogous to CLOS's
\texttt{defclass}\index{defclass@\texttt{defclass}}.

For layered applications and embedded languages, CL provides an environment unmatched
by other current languages. Because CL allows macros to generate other macros, and the
Lisp code can compute and generate Lisp code at all levels, it is straightforward to build
extremely high-level languages on top of CL. And because of the CL-based syntax, these
layered languages integrate seamlessly into the underlying CL environment.

\section{Conclusion}

Chapter 4 has provided a brief introduction into some of the interfacing capabilities available
in CL. For additional information on any of these topics, you should refer to the on-line
and printed documentation from Franz Inc., as well as to various resources available on the
Web.
